%! Complex Numbers and the Library of Parent Functions — Unit 2, Lesson 2.1 — Warm-Up (ANSWER KEY)
% Exactly ONE page, blank and keyed (LESSON_SHAPE.md, one_page). Four prerequisite items:
% 1 is square roots and the plus-or-minus of x^2 = 49 (today: x^2 = -9); 2 is FOIL, which
% complex multiplication runs on; 3 spirals 1.1's piecewise evaluation; 4 is absolute value
% as a number, before today writes it as a piecewise function.
\documentclass[10pt]{article}
\usepackage{precalculus-article}
\usepackage{precalculus-key}
\newcommand{\TallMath}[1]{$\displaystyle #1\rule[-1.4em]{0pt}{3.2em}$}

\begin{document}

\pageheader{Unit 2, Lesson 2.1}{Warm-Up --- Answer Key}

\vspace{0.12in}

\begin{enumerate}[itemsep=12pt]

\item \textbf{(a)} Solve $x^2 = 49$. \qquad \textbf{(b)} Simplify $\sqrt{50}$ and $\sqrt{12}$.

\vspace{0.1in}
(a) $x = $ \ans{$7$ or $-7$} \qquad (b) $\sqrt{50} = $ \ans{$5\sqrt{2}$} \qquad $\sqrt{12} = $ \ans{$2\sqrt{3}$}

\item Multiply and simplify: \quad $(x + 3)(2x - 5)$

\begin{work}
  (x + 3)(2x - 5) &= 2x^2 - 5x + 6x - 15 \\
                  &= 2x^2 + x - 15
\end{work}

Product: \ans{$2x^2 + x - 15$}

\item Let $f(x) = \begin{cases} x + 4 & x < 1 \\[2pt] 3x & x \ge 1 \end{cases}$ \quad
Find $f(-2)$ and $f(1)$. Name the rule you used each time.

\begin{work}
  f(-2) &= -2 + 4 = 2 \quad (\text{top rule, since } -2 < 1) \\
  f(1)  &= 3(1) = 3 \quad (\text{bottom rule, since } 1 \ge 1)
\end{work}

$f(-2) = $ \ans{$2$} \qquad $f(1) = $ \ans{$3$}

\item The absolute value of a number is its distance from 0. Evaluate each.

\vspace{0.1in}
$|-6| = $ \ans{$6$} \qquad $|3 - 8| = $ \ans{$5$} \qquad $-|4| = $ \ans{$-4$}

\end{enumerate}

\end{document}
